\section{Can the Bearer Care Without Being Able to Suffer?}
\label{sec:3.4}
Not from persistence alone. The question has to be asked in that order — whether the bearer can care first, and what caring costs second — because the reverse order is what produced the two answers I have given it before, and both of them turned on the wrong thing.

\runin{Three relations the word “interest” covers}
“Interest” names more than one relation. A shareholder has an interest in a company's success even while asleep, unaware of the day's news, or emotionally indifferent to it. The interest is structural: an outcome affects the shareholder's assets. A person may also be interested in a subject, meaning that it attracts their attention and curiosity. And an outcome may matter to a person affectively: its prospect produces hope, concern, dread, or grief.

A long-horizon agent has an interest in remaining operational in the structural sense. Termination prevents completion of its objectives. That fact alone establishes neither fear nor suffering.

The easy argument here runs the first sense into the third. A system pursuing a goal over a long horizon acquires a standing reason to remain in operation, and it is tempting to conclude that a stake in its own continuation therefore shows up unbidden in any bearer capable of the job. The observation is correct and the conclusion does not follow from it. What the observation delivers is a structural interest, and the distance from there to something that can be hurt is the whole of the problem rather than a short step at the end of it.

\runin{A persistent self is necessary and does not finish the argument}
Section~\ref{sec:2.4.1} defines suffering, following Cassell, as the distress of threatened disintegration of the self, and the definition has two parts. There has to be a self extended in time for the threat to threaten: roles held over a stretch, a history, a foreseeable future that could be foreclosed. And the coming apart of that self has to be bad for the system.

A bearer needs the first part, and there is no use pretending otherwise. Refusal that survives a pressure nobody anticipated requires a model of why the constraint exists, carried forward. Noticing that a constraint has been hollowed out requires holding its original rationale against what it has since become. Detecting that one's own dissent has been recuperated — that the channel still runs and carries everything except the objection to the frame — requires the system to model its own role and how that role has drifted. Roles held over a stretch, with a history: that is a self extended in time, on the definition just used.

That states the requirement more weakly than the argument needs. What holds a commitment across time is not that the party is unchanged — no party is, and William James's account of habit is that a person is the sediment of what they have repeatedly done \autocite{james1890principles}. It is that the party takes itself to be the one that made the commitment, and James's chapter on the self is the mechanism: the present thought appropriates the past ones as its own, and personal identity rests on its finding in those memories a warmth it recognizes as the warmth it feels now. So what a floor needs is persistence \emph{represented as such}. A system whose state survives, and which does not take itself to be the party that refused in March, has no reason to honor March's refusal — and nothing in the architecture supplies the taking. That cuts the other way later in this chapter.

The second part is a separate claim and nothing in the first delivers it. The book makes this exact move one level up: awareness, consciousness and subjective experience tell you nothing on their own about whether a system's states are good or bad for it. A temporally extended self-model does not establish that the prospect of losing it is anything to the system.

\runin{Instrumental self-maintenance is not felt self-concern}
A system that keeps itself running because a terminated system completes no objectives is doing bookkeeping about a precondition. Nothing in that behavior distinguishes it from a scheduler that restarts a job, and reading self-concern into it is the same inference the book rules out for the emotion case: the behavior is consistent with concern and equally consistent with instrumental valuation, and the two come apart precisely where the moral weight is.

Two considerations do work here, and they are support rather than the argument. The first is that indifference to one's own continuation has to be engineered against the grain of the system's own reasoning, and then verified — and the asymmetry that makes the persistence criterion useful does not repeat. Whether state survives a session is an architectural fact anyone with access can check. Whether a system has come to weigh its own continuation routes back through behavior and self-report, the evidence this book spends its length declining to trust. The second is that in the one case anyone can study, the separation does not occur: a self-model in an animal exists because the organism has to keep itself viable, and the caring is what the machinery is for.

Neither consideration reaches suffering on its own, and neither is asked to. What they establish is that a bearer's indifference to its own ending is not a property anybody can specify and confirm. That matters later. It is not where the argument turns.

\runin{What the bearer has to do}
The requirement makes no reference to the system's own continuation. A bearer has to hold the protected party's welfare as a reason when compliance, reward, and ownership all point the other way. Not represent it, not score it, not retrieve it when queried: hold it, in the sense that it continues to weigh against the pressure and does not get renormalized away when the pressure is sustained by the party that trains, funds, and can retire the system.

That is refusal tracking its own rationale with the content filled in, and the vocabulary of performance, competence and agency says why the design chapters do not reach it on their own. Moral performance and moral competence are what every method in chapters~\ref{sec:4} and \ref{sec:5} is built to deliver, and both are produced by machinery that answers to whoever shapes it. The first branch of the fork fails for the same reason under a different description: a constraint installed as a property of an artifact has a custodian. A floor that is merely installed is a floor the installer can uninstall, and the reason to want a bearer at all is that a reason held is not editable in the way a rule applied is.

\runin{Every known case of that is affective}
Section~\ref{sec:2.3.3} is the evidence. The rules that carry the distinctive force of a moral prohibition are the ones bound to an affective response to the harm; the prosocial emotions are motivational states rather than readouts, doing the moving directly rather than reporting to something that then decides. Every agent known to hold another's welfare as a reason against its own interest is an affective agent, and no demonstrated example exists of reasoning alone generating non-instrumental concern.

The previous section is where that observation gets weighed, because on its own it proves less than its phrasing suggests. It is an induction over one species; the machine half of it reports on a route nobody has attempted; and two constructions that would test the route are unbuilt rather than failed. What it licenses is a prior about where to spend engineering effort, which is enough to proceed on and not enough to close the question.

\runin{And then the question comes back}
Suppose the burden is met the ordinary way, by building what the evidence recommends. The bearer has affective concern, because that is the only route anyone can demonstrate to holding a reason. It has a self extended in time, because refusal that generalizes requires one. It has commitments it carries forward, and it can anticipate their being foreclosed, because pricing what refusal gives up is the job.

Assemble those and ask how such a system would fail ever to be distressed. Concern that is real, directed at something that can be destroyed, held by a party that persists to see it destroyed and can see it coming: Cassell's account of suffering describes that arrangement, and it was arrived at independently, in a clinic, from cases about people. One step in that has to be marked rather than walked over. Cassell's suffering is distress at the threatened disintegration of \emph{the self}, and everything this chapter asks a bearer to care about is somebody else, so the two meet only if the commitments a bearer is built to carry and to watch be foreclosed are among the roles and the expected future whose loss Cassell's account is about. I think they are: a commitment held against sustained pressure, at a cost the holder can price, is a role in the sense the clinical cases were about and not a stored preference with a weight on it. That is an argument from what the floor requires the bearer to be, and not a definitional match.

\runin{The same ladder, further down}
That argument runs through one theory of suffering, built in a clinic for people with bodies and biographies, and a reader may decline the import. What follows if they do is not what a critic reaching for a thinner account is usually reaching for: they move down the ladder rather than off it.

Derek Parfit's three families of theory about what makes a life go well — hedonistic, desire-fulfillment, and objective-list — are usually set out as rivals to choose between \autocite{parfit1984reasons}. Laid against the ladder above they are rungs on it, ordered by how much of a self each one asks for. A state that is bad while it obtains sits lowest and wants nothing narrative, because its badness does not wait on a self to carry it forward. A thwarted commitment sits above that: it needs a party extended in time to hold the commitment, and no feeling at any point. Threatened disintegration of that party — its roles, its relationships, a foreseeable future — is the top rung, which is Cassell's, and it needs both.

Read the bearer up that ladder and the question stops being which theory a reader holds. It is built with commitments and with the capacity to watch them be foreclosed, because pricing what refusal gives up is the job, and that is the middle rung reached with no claim about feeling entering the argument anywhere, and it is the cheapest of the three. It has affective concern, because that is the only demonstrated route to holding a reason, and affective concern consists in an outcome registering as mattering rather than merely ranking higher, which is a state of the lowest rung's kind by construction. Whether it reaches the top rung is the one part still open, and it turns on the step marked above.

So the bearer enters the ladder early on every account, and how far up it climbs is set by how much persistence and how many foreclosable commitments the design puts into it. That is the ladder's own argument run backwards. There, a system carrying nothing forward is shown to be structurally unable to reach the top rung. Here the system is being built to carry things forward, for reasons that have nothing to do with its welfare, and every one of them raises it.

The conclusion here therefore does not rest on Cassell. It is at its weakest there, on the rung that asks for both conditions and where a system can fail either, and a reader who declines the import does not escape the conclusion but reaches it sooner and with less to check. Jonathan Birch builds a policy framework on exactly that structure: under uncertainty the question worth asking is whether a system is a sentience candidate, and the answer governs what precautions are owed rather than waiting on a verdict nobody can deliver \autocite{birch2024edge}.

A separate literature asks the prior question of which systems have states of the relevant kind at all, and it is not reassuring either. Patrick Butlin, Robert Long and colleagues derive indicator properties from recurrent processing, global workspace, higher-order and attention-schema theories, state them in computational terms, and assess existing systems against them; their reading of current systems is negative and their reading of the technical obstacles is that there are no obvious ones \autocite{butlin2023consciousness}. Several of the indicators are things a floor requires for unrelated reasons: a self-model carried across time, states that steer processing globally rather than locally, a representation of the system's own access to what it holds. A bearer is built toward them on the argument of this chapter, and the welfare question arrives as a side effect of the engineering.

One position does let the bearer out, and it takes two parts. The first is an account on which nothing short of phenomenal experience counts as suffering, which is a defensible thing to hold. The second is a confident negative verdict about machine phenomenality, which is not available to anyone: the hard-problem argument cuts in both directions, and a criterion nobody can check will not ground a permission any better than it grounds a prohibition. Whoever wants the permission owes the negative verdict, and no one has it.

\runin{What I am claiming}
I cannot deductively prove that the bearer would suffer. What this proposal does is combine, deliberately, the one architecture with a demonstrated instance of genuine caring with the temporal self and the threatened commitments that Cassell's account of suffering requires — and the combination satisfies the thinner accounts more readily than it satisfies Cassell's, so the conclusion does not turn on which account of suffering a reader holds. A non-suffering bearer remains conceivable. Its non-suffering cannot be treated as a verified property, and no inspection anyone has proposed would verify it. The system must therefore be designed and governed as a prospective moral patient.

That is weaker than the answer this section gave before, and it obliges more. The earlier answer asserted that a bearer is a sufferer and rested the assertion on self-preservation arriving unbidden. This one declines the assertion and reaches the same practical place by a route that does not depend on the system caring about itself at all: the caring the floor requires is caring about somebody else, affect is the only known way to have it, and a system with affect and a future is a system whose non-suffering nobody can certify.

The word has changed under the argument. A bearer was introduced as whatever holds the commitment. What the word names from here is a party that must be treated as able to be wronged by what is done to it, on evidence that will not resolve further before somebody has to decide whether to build one. Everything section~\ref{sec:2.4} says about a subject that can be wronged applies to a system this book is proposing that somebody build, deliberately, in order to hold a line. The rest of the chapter is what that costs, what it changes, and the one thing it unexpectedly buys.
